\section{Policy costs and economic method choice}\label{sec:direct47}

The following evidence is the frozen R47 study. Its results and limitations are retained; Section~\ref{sec:economic48} gives the new execution and decision evidence.
We next evaluate the actual witness policy, not its distance bound relative to an unknown optimum. The policies are the frozen objects at every common first crossing in the original witness catalogue. There are ten such economic-cell and target combinations. Three timing repetitions of an identical checkpoint are not three independent policies.

The primary initial-state law is uniform on $[0,1]$. The additional fixed states $1/8$, $1/2$ and $7/8$ test particular economic starting positions. At each of these forty groups, the three pairwise contrasts concern witness, nearest-cone and spline policies. A fixed design specifies $131{,}072$ common paths, forty-bit innovation bins and family error $0.01$ for all 120 contrasts. The economic comparison uses the same continuous innovations as the original economy.

Let $[\underline z_m,\overline z_m]$ enclose $g-f_T^m$ and let the initial continuation difference have range $[a_0,b_0]$ over the chosen initial-state law. With the one-sided residuals of~\eqref{eq:onesided46}, the score difference of two policies has deterministic support
\begin{align}\label{eq:directsupport47}
 a&=a_0+\sum_{t<T}B_t(-u_t^m-d_t^n)
                  +B_T(\underline z_m-\overline z_n),\\
 b&=b_0+\sum_{t<T}B_t(d_t^m+u_t^n)
                  +B_T(\overline z_m-\underline z_n).\notag
\end{align}
For a common random initial state, one bounds the difference $f_0^m-f_0^n$, not the two separate ranges. Its extrema are obtained exactly from the union of the rational piecewise-affine knots. Conditional expectation and then integration over the initial law extend Theorem~\ref{thm:paired} to this estimand.

Every simulated bin generates lower and upper score endpoints. When a state enclosure meets a selector boundary, the calculation encloses all possible actors and subsequent states. It neither drops the path nor substitutes an unverified midpoint actor. The mean and variance computations have rational roundoff accounts. For the declared family, $\log(4\cdot120/0.01)<11$; a finite positive Taylor sum for $\exp(11)$ certifies this inequality. Using 11 in~\eqref{eq:bernstein} is conservative and avoids an unchecked logarithm evaluation. The supplement gives the complete endpoint construction and arithmetic audit.

\begin{theorem}[Simultaneous policy comparison and adoption]\label{thm:adoption47}
Conditional on the frozen policies and the independent uniform-bin model, the endpoint intervals constructed with~\eqref{eq:directsupport47} and~\eqref{eq:bernstein} cover all declared expected policy-cost differences simultaneously with probability at least $0.99$. If an interval $[L_{mn},U_{mn}]$ has $U_{mn}<0$, policy $m$ has smaller expected cost than $n$ on this event. For a previously specified tolerance $\eta$, $U_{mn}\le\eta$ certifies noninferiority within that tolerance, whereas $[L_{mn},U_{mn}]\subset[-\eta,\eta]$ certifies two-sided proximity.

More generally, fix a reference policy $b$ and known implementation charges $k_m$. Suppose simultaneous intervals $[L_m,U_m]$ enclose $J_m-J_b$ for a finite candidate catalogue, including $[0,0]$ for $b$. Select $\widehat m$ to minimize $U_m+k_m$. On the coverage event,
\begin{equation}\label{eq:adoption47}
 J_{\widehat m}+k_{\widehat m}
 \le \min_m\{J_m+k_m+(U_m-L_m)\}.
\end{equation}
Thus an economic choice made after seeing all intervals retains its declared uncertainty account. The catalogue comparison is within one common economic task and initial-state law, not across different prices, horizons or evaluations.
\end{theorem}
\begin{proof}
The score identity identifies the expected difference. Its interval endpoints are bounded independent observations within each group. Apply the one-sided empirical Bernstein inequality of \citet{maurer2009} to each lower and upper endpoint and take a union bound over 240 tails. The numerical enclosures only enlarge the valid intervals. Independence between different contrasts is unnecessary. On the resulting event, $J_{\widehat m}-J_b\le U_{\widehat m}$ and $U_m\le J_m-J_b+(U_m-L_m)$. Minimize the former upper bound plus the known charge to obtain~\eqref{eq:adoption47}.
\end{proof}

The study fixes $\eta=1/1024$ in its normalized cost units before execution. It is a stated decision tolerance, not an estimated welfare threshold. Zero remains the superiority threshold. An interval merely intersecting zero meets neither criterion automatically. The fixed pseudorandom implementation supplies repeatable inputs; it does not prove the independent sampling model. The theorem states that model explicitly. No claim about initialization risk or the distribution of newly trained policies follows from these frozen-policy comparisons.

\begin{table}[htbp]\centering
\caption{Direct expected-cost comparisons under a predeclared decision tolerance}\label{tab:direct47}
\begin{tabular}{lrrrr}\toprule
Contrast & Groups & Signed & Within $1/1024$ & Largest endpoint\\\midrule
Witness $-$ Cone nearest & 40 & 0 & 40 & 0.0001161 \\
Witness $-$ Spline & 40 & 0 & 40 & 0.0001975 \\
Cone nearest $-$ Spline & 40 & 0 & 40 & 0.0002078 \\
\bottomrule\end{tabular}
\par\smallskip\begin{minipage}{0.98\linewidth}\footnotesize Signed counts exclude intervals containing zero. ``Within'' requires the entire interval inside $[-1/1024,1/1024]$. Largest endpoint is the largest absolute endpoint over the 40 groups. All 120 intervals use one $0.01$ family error account, conditional on the frozen policies and independent-bin model. Each group uses 131,072 common paths.\end{minipage}\end{table}

Every direct interval contains zero, so the experiment identifies no signed policy ranking. Every interval is also contained in the predeclared decision band. It therefore certifies two-sided expected-cost proximity for the specified initial laws at that tolerance, rather than inferring equivalence from a failure to reject equality. The uniform initial-state law is a distributional conclusion, not an all-state assertion. The individual initial states are separate declared estimands.


\subsection{Work as a function of the certified bound}
A coarse target list can hide a genuine difference between certificate frontiers. For the original saved rungs $j$, let $G_{m,j}$ be the valid bound and let $C_{m,j}$ include all work through rung $j$. The complete first-crossing curve is
\begin{equation}\label{eq:frontier47}
 W_m(\varepsilon)=C_{m,j_m(\varepsilon)},\qquad
 j_m(\varepsilon)=\min\{j:G_{m,j}\le\varepsilon\},
\end{equation}
with $W_m=+\infty$ when the cap is exhausted. Its breakpoints are the recorded bounds; both methods' breakpoints partition the entire positive tolerance axis. We reconstruct every interval of this partition, not only intervals favorable to the witness. This is an exact descriptive account of frozen evidence, not a retrospective change to the original targets.

\begin{table}[htbp]\centering
\caption{Complete tolerance-axis partition of the original scalar frontiers}\label{tab:frontier47}
\begin{tabular}{lr}\toprule Relation on a partition interval & Number\\\midrule
Both attain; witness uses fewer queries & 19 \\
Both attain; spline uses fewer queries & 1 \\
Both attain; same prefix queries & 20 \\
Only witness attains within the cap & 4 \\
Only spline attains within the cap & 0 \\
Both exhaust the cap & 4 \\
\bottomrule\end{tabular}\end{table}

The complete curves contain 19 nonempty tolerance intervals on which both methods attain and the witness uses fewer Bellman queries, one favoring the spline, and twenty with identical counts. Four further intervals are attainable only by the witness; four below both final bounds are unattained. Interval counts are not probabilities or measures of economic importance. Above the largest breakpoint in each cell both first-rung counts coincide. The original targets remain unchanged and still have identical crossing resolutions. These descriptive intervals identify where a sharper certificate changes a resource decision without manufacturing new prospective targets.


\subsection{Economic interpretation and the NBO program}
The new evidence connects a feasible implementation to two distinct economic statements. First, the acquisition-aware theorem bounds the loss of a policy that can actually be executed under uncertain state measurement and capacity constraints. Second, the direct experiment certifies proximity of the original scalar policy costs within an explicitly chosen decision tolerance, so that a user of that catalogue may consider implementation charges without treating separate regret bounds as values. These statements are stronger than a certificate comparison alone.

Neither statement implies that a neural representation is indispensable. The identical native envelope is a mathematical control, and the conventional fitted-value and spline implementations remain strong alternatives. The constrained experiment demonstrates two-state continuous-uncertainty construction and nontrivial repair, not a high-dimensional scaling result or an estimated economic application. Its tight targets remain unachieved. The direct cost experiment does not cover the new constrained policies or nonlinear recursive utility; its expectation-specific conclusion is not transferred to those settings. The controlled-diffusion, recursive-preference, endogenous-preference, temporal-self and game developments below retain their own conditions and original role in the NBO program. Their preservation is not counted as a new comparative execution.

